\documentclass[11pt]{article}

\usepackage[margin=1in]{geometry}
\usepackage{amsmath,amssymb,amsfonts,amsthm}
\usepackage{array}
\usepackage{bm}
\usepackage{graphicx}
\usepackage{multirow}
\usepackage{multicol}
\usepackage{nicefrac}
\usepackage{paralist}
\usepackage{enumitem}
\usepackage{verbatim}
\usepackage{thm-restate}
\usepackage[normalem]{ulem}
\usepackage[compact]{titlesec}
\usepackage{tikz}
\usetikzlibrary{arrows}
\usetikzlibrary{arrows.meta}
\usetikzlibrary{backgrounds}
\usetikzlibrary{calc}
\usetikzlibrary{decorations.markings}
\usetikzlibrary{fit}
\usetikzlibrary{patterns}
\usetikzlibrary{positioning}
\usetikzlibrary{shapes}

\usepackage{fancybox}
\newenvironment{fminipage}%
  {\begin{Sbox}\begin{minipage}}%
  {\end{minipage}\end{Sbox}\fbox{\TheSbox}}

\newenvironment{algbox}[0]{%
\vskip 0.2in
\noindent
\begin{fminipage}{6.3in}
}{%
\end{fminipage}
\vskip 0.2in
}

\usepackage{tcolorbox}
\tcbuselibrary{skins,breakable}
\tcbset{enhanced jigsaw}
\usepackage{mdframed}
\usepackage{xcolor}
\usepackage{cite}
\usepackage{url}
\usepackage{xspace}
\usepackage[mathscr]{euscript}
\usepackage{nameref}
\usepackage[
  linktocpage=true,
  pagebackref=true,
  colorlinks=true,
  linkcolor=blue,
  citecolor=blue,
  urlcolor=blue,
  bookmarks,
  bookmarksopen,
  bookmarksnumbered
]{hyperref}
\usepackage[noabbrev,nameinlink]{cleveref}

\newtheorem{theorem}{Theorem}[section]
\newtheorem{example}{Example}[section]
\newtheorem{lemma}{Lemma}[section]
\newtheorem{proposition}[lemma]{Proposition}
\newtheorem{corollary}[lemma]{Corollary}
\newtheorem{claim}[lemma]{Claim}
\newtheorem{fact}[lemma]{Fact}
\newtheorem{conjecture}[lemma]{Conjecture}
\newtheorem{definition}[lemma]{Definition}
\newtheorem{problem}{Problem}
\newtheorem{remark}[lemma]{Remark}
\newtheorem{observation}[lemma]{Observation}

\newtheorem*{theorem*}{Theorem}
\newtheorem*{claim*}{Claim}
\newtheorem*{proposition*}{Proposition}
\newtheorem*{lemma*}{Lemma}
\newtheorem*{problem*}{Problem}

\allowdisplaybreaks
\setlength{\parskip}{3pt}

\title{A Diagonalizable Obstruction for Random Two-Step Ritz Compressions}
\author{}
\date{August 3, 2026}

\begin{document}

\pagenumbering{roman}

\maketitle

\begin{abstract}
We consider two-step Krylov compression of an arbitrary real matrix
$A\in\mathbb R^{n\times n}$ from a standard Gaussian starting vector
$b\sim N(0,I_n)$.  On the event that $\mathcal K_2(A,b)$ is two-dimensional,
let $Q$ denote any orthonormal basis of this space; the associated Ritz matrix
is $Q^TAQ$.  For every even
$n\geq4$, we construct a real diagonalizable matrix $A_n$ with
$\lVert A_n\rVert_2\leq1+e^{-n}/2<3/2$ such that
$\mathcal K_2(A_n,b)$ is two-dimensional
almost surely and
\[
  \Pr\!\left\{\kappa_V(Q^TA_nQ)\geq \frac{e^n}{5}\right\}
  \geq 1-2e^{-n/40}.
\]
For $\varepsilon_n=5e^{-2n}$, the same family satisfies
\[
  \mathbb E\!\left[\operatorname{area}
  \Lambda_{\varepsilon_n}(Q^TA_nQ)\right]
  \geq \left(1-2e^{-n/40}\right)\frac{\pi e^{-2n}}4.
\]
Consequently, for every positive polynomial $p$, the inequality
$\kappa_V(Q^TA_nQ)\leq p(n)$ fails with probability tending to one
exponentially fast along the even dimensions.  Moreover, for every fixed
$\beta>1$ and every positive polynomial $p$, no upper bound of the form
$p(n)\varepsilon^\beta$ holds uniformly over all dimensions, inputs, and
positive tolerances under this Gaussian starting-vector model.
\end{abstract}

\noindent\textbf{Note.} This paper was fully AI generated by an internal
harness using GPT 5.6 Sol xhigh, with a minimal prompt from Richard Peng.

\clearpage
\tableofcontents
\clearpage

\pagenumbering{arabic}

\section{Introduction}
\label{sec:introduction}

Krylov subspace methods replace a large matrix eigenvalue problem by a
projection onto a space generated from a starting vector.  Given
$A\in\mathbb C^{n\times n}$ and $b\in\mathbb C^n$, the $k$-step space and its
orthogonal compression are
\[
  \mathcal K_k(A,b)=\operatorname{span}\{b,Ab,\ldots,A^{k-1}b\},
  \qquad H_k=Q^*AQ,
\]
where the columns of $Q$ form an orthonormal basis for
$\mathcal K_k(A,b)$.  The eigenvalues of the smaller matrix $H_k$ are the Ritz
values produced by the Arnoldi process, a basic tool for large nonsymmetric
eigenvalue problems~\cite{Arnoldi1951MinimizedIterations,Saad1980VariationsArnoldi}.
For a nonnormal compression, however, the location of the Ritz values is only
part of the stability question.  The condition number of an eigenvector
matrix governs eigenvalue sensitivity through the Bauer--Fike theorem, while
the pseudospectrum records the same sensitivity without requiring
diagonalizability~\cite{BauerFike1960NormsExclusion,Trefethen1999ComputationPseudospectra}.

Randomizing the starting vector is a natural way to avoid exceptional Krylov
spaces.  Such starts have long supported positive convergence guarantees for
the power and Lanczos methods on symmetric positive-definite
matrices~\cite{KuczynskiWozniakowski1992RandomStart}.  That setting does not,
though, test nonnormal eigenvector stability: if $A$ is Hermitian, then every
$Q^*AQ$ is Hermitian, has eigenvector condition number one, and has
$\varepsilon$-pseudospectral area at most $\pi k\varepsilon^2$.  At the other
extreme, deterministic Arnoldi compressions have considerable adversarial
freedom.  Duintjer Tebbens and Meurant showed that prescribed Ritz-value
histories can be realized~\cite{DuintjerTebbensMeurant2012AnyRitz}; this is a
result about Ritz values for chosen matrix--vector pairs, rather than about
eigenvector conditioning under a random start.

Amsel et al.~\cite{AmselEtAl2026LinearSystems} posed the corresponding
random-start stability question for an arbitrary matrix $A$.  They asked
whether $\kappa_V(H_k)$ is bounded by a polynomial in $n$ with high
probability, or alternatively whether the expected pseudospectral area admits
a bound of the form
\[
  \mathbb E\!\left[\operatorname{area}\Lambda_\varepsilon(H_k)\right]
  \leq \operatorname{poly}(n)\varepsilon^\beta
\]
for an exponent $\beta$ close to $2$.  Their formulation also records an
important deterministic warning: for the cyclic shift and the special start
$b=e_1$, the compression before full dimension is a Jordan block.  The point
of randomizing $b$ is precisely to ask whether such badly aligned starts are
negligible.  The distribution of the random vector and the intended range of
$\varepsilon$ are not fixed in that formulation.  Here we use the standard
real-Gaussian model $b\sim N(0,I_n)$ and consider the literal quantification
over arbitrary inputs.

The closest positive comparison changes the random object.
Shah~\cite{Shah2025RandomCompressions} studies the same expected-area observable
when $Q$ spans an independent Haar-random complex subspace, proving bounds
with $\varepsilon$-exponents $6/5$, $4/3$, and $2$ (up to polynomial and
logarithmic factors) under successively stronger hypotheses on $A$.  A
random-start Krylov space is not a Haar-random plane: its distribution is
coupled to $A$ through $b,Ab,\ldots$.  Our result shows that this distinction
is decisive for arbitrary inputs.  Even at $k=2$, randomizing the start does
not prevent exponentially ill-conditioned Ritz eigenvectors, and the same
examples rule out every uniform expected-area exponent greater than one.
The precise statement is as follows.

\begin{theorem}[Diagonalizable obstruction for two-step Ritz compressions]
\label{thm:main}
For every even $n\geq4$, there exists a diagonalizable matrix
$A_n\in\mathbb R^{n\times n}$ with
$\lVert A_n\rVert_2\leq1+e^{-n}/2<3/2$ having the
following properties.  Let $b\sim N(0,I_n)$ and set
\[
  \mathcal K_2(A_n,b)=\operatorname{span}\{b,A_nb\}.
\]
This Krylov space is two-dimensional almost surely.  On this event, let $Q$
be any orthonormal basis of the Krylov space.  Then
\[
  \Pr\!\left\{\kappa_V(Q^TA_nQ)\geq \frac{e^n}{5}\right\}
  \geq 1-2e^{-n/40}.
\]
Moreover, for $\varepsilon_n=5e^{-2n}$,
\[
  \mathbb E\!\left[\operatorname{area}
  \Lambda_{\varepsilon_n}(Q^TA_nQ)\right]
  \geq \left(1-2e^{-n/40}\right)\frac{\pi e^{-2n}}4.
\]
Consequently, no positive polynomial $p(n)$ bounds $\kappa_V(Q^TA_nQ)$ with
high probability uniformly over arbitrary inputs, and for every fixed
$\beta>1$ no upper bound of the form $p(n)\varepsilon^\beta$ holds uniformly
for the expected pseudospectral area.
\end{theorem}

Several features of the theorem delimit the obstruction.  First, it is not a
defective-input example: $A_n$ is diagonalizable and uniformly bounded in
operator norm.  Its two eigenspaces are nevertheless exponentially close.
The minimal polynomial has degree two, so the two-step Krylov plane is almost
surely invariant; with probability $1-e^{-\Omega(n)}$, the two spectral
components of the Gaussian start are nearly antiparallel.  Orthogonal
compression preserves their angle, which becomes the angle between the two
Ritz eigenvectors.  A direct two-dimensional Schur-form estimate then places
a disk of radius $e^{-n}/2$ in the $5e^{-2n}$-pseudospectrum.  Thus the
conditioning and area conclusions arise from the same geometry.

Second, the expected-area conclusion is witnessed at the exponentially small
tolerance $\varepsilon_n=5e^{-2n}$.  It therefore rules out estimates intended
to hold uniformly for all positive $\varepsilon$, but does not rule out a
statement restricted to inverse-polynomial tolerances.  Likewise, the
construction does not address normal inputs, the cyclic shift, prescribed
growing Krylov dimensions, a complex-Gaussian start, or bounds allowed to
depend on the conditioning of an eigenbasis for $A$.  These are genuine
structural distinctions.  For example, unitary Arnoldi admits an
orthogonal-polynomial description~\cite{Gragg1993ArnoldiIsometric}, while
known results on Ritz values of normal matrices impose geometry absent from
general nonnormal inputs~\cite{CardenHansen2013NormalRitz}.  Neither line is a
conditioning theorem for the random-start Krylov compression considered
here.

Since nonzero rescaling of $b$ does not change its Krylov space, our Gaussian
model is equivalently a uniform random direction on the real unit sphere; it
should not be confused with an independent Haar-random two-dimensional
subspace.  \Cref{sec:preliminaries} fixes the conditioning and pseudospectral
conventions.  \Cref{sec:overview} explains the construction and its geometric
mechanism, and \Cref{sec:proof} proves the main theorem.

\section{Preliminaries}
\label{sec:preliminaries}

Throughout, vector norms are Euclidean and matrix norms are the induced
operator norms.  A standard Gaussian vector $b\sim N(0,I_n)$ has independent
$N(0,1)$ coordinates.  In particular, if $n=2r$ and $b=(u,v)$ under the block
decomposition $\mathbb R^n=\mathbb R^r\oplus\mathbb R^r$, then $u$ and $v$ are
independent $N(0,I_r)$ vectors.

\paragraph{Krylov spaces and Ritz matrices.}
For $A\in\mathbb R^{n\times n}$, $b\in\mathbb R^n$, and $k\geq1$, define
\[
  \mathcal K_k(A,b)
  =\operatorname{span}\{b,Ab,\ldots,A^{k-1}b\}.
\]
If this space has dimension $d$ and
$Q\in\mathbb R^{n\times d}$ has orthonormal columns spanning it, then
$H=Q^TAQ$ is the associated \emph{Ritz matrix}.  A different orthonormal basis
replaces $H$ by an orthogonally similar matrix.  If $\mathcal K_k(A,b)$ is
$A$-invariant, then
\begin{equation}
  AQ=QH.
  \label{eq:invariant-compression}
\end{equation}
In particular, if $x\in\mathcal K_k(A,b)$ is an eigenvector of $A$, then
$Q^Tx$ is an eigenvector of $H$ with the same eigenvalue.  Moreover, $Q^T$
preserves norms and inner products on $\mathcal K_k(A,b)$.

\paragraph{Angles and eigenvector conditioning.}
For nonzero $x,y\in\mathbb R^d$, let
\[
  c(x,y)=\frac{|x^Ty|}{\lVert x\rVert_2\lVert y\rVert_2},
  \qquad
  s(x,y)=\sqrt{1-c(x,y)^2}.
\]
Thus $s(x,y)$ is the sine of the acute angle between the lines spanned by
$x$ and $y$, and
\begin{equation}
  s(x,y)^2
  =\frac{\lVert x\rVert_2^2\lVert y\rVert_2^2-(x^Ty)^2}
  {\lVert x\rVert_2^2\lVert y\rVert_2^2}.
  \label{eq:gram-angle}
\end{equation}
If a real diagonalizable $2\times2$ matrix $H$ has distinct eigenvalues, let
$V=[x_+\;x_-]$ have unit eigenvector columns and define
$\kappa_V(H)=\lVert V\rVert_2\lVert V^{-1}\rVert_2$.  This definition is
unchanged by signs or ordering of the two eigenvectors.  Writing
$c=c(x_+,x_-)$ and $s=s(x_+,x_-)$, the eigenvalues of $V^TV$ are $1+c$ and
$1-c$, so
\begin{equation}
  \kappa_V(H)
  =\sqrt{\frac{1+c}{1-c}}
  =\frac{1+c}{s}
  \geq\frac1s.
  \label{eq:angle-conditioning}
\end{equation}
Orthogonal similarity preserves this quantity.

\paragraph{Pseudospectra.}
For a square matrix $H$, write $\sigma_{\min}(H)$ and $\sigma_{\max}(H)$ for
its smallest and largest singular values, respectively, and define its
$\varepsilon$-pseudospectrum by
\[
  \Lambda_\varepsilon(H)
  =\{z\in\mathbb C:\sigma_{\min}(zI-H)\leq\varepsilon\}.
\]
Area means planar Lebesgue measure on $\mathbb C$.  Orthogonal similarity
preserves singular values and hence preserves the pseudospectrum.  Finally,
for every $2\times2$ matrix $M$,
\begin{equation}
  \sigma_{\min}(M)\sigma_{\max}(M)=|\det M|.
  \label{eq:two-by-two-singular-values}
\end{equation}
We also use the elementary bound $\sigma_{\max}(M)=\lVert M\rVert_2\geq
|M_{ij}|$ for each entry of $M$.

\section{Overview}
\label{sec:overview}

The proof is a two-eigenvalue counterexample whose geometry survives Krylov
compression exactly.  Write $n=2r$, set
$\delta=e^{-n}$ and $\alpha=\delta/2$, and consider
\begin{equation}
  A_n=
  \begin{pmatrix}
    \alpha I_r&-I_r\\
    0&-\alpha I_r
  \end{pmatrix}.
  \label{eq:overview-input}
\end{equation}
Its eigenspaces are
$E_+=\{(x,0):x\in\mathbb R^r\}$ and
$E_-=\{(x,\delta x):x\in\mathbb R^r\}$, with eigenvalues
$\alpha$ and $-\alpha$, respectively.  Thus the eigenspaces are separated by
an angle of order $\delta$, but $A_n$ is nevertheless diagonalizable and has
$\lVert A_n\rVert_2\leq1+\alpha<3/2$.  The choice
$2\alpha=\delta$ is what reconciles these two features: the exponentially
ill-conditioned eigenbasis is multiplied by an equally small eigenvalue gap,
leaving the off-diagonal block of $A_n$ at constant scale.

We take the starting vector to be $b\sim N(0,I_n)$.  For this compression the
model is equivalent to choosing a uniform random direction, because rescaling
$b$ does not change its Krylov space.  Writing $b=(u,v)$, split it into its two
spectral components
\begin{equation}
  p=(u-\delta^{-1}v,0),
  \qquad
  m=(\delta^{-1}v,v).
  \label{eq:overview-components}
\end{equation}
They satisfy $b=p+m$, $A_np=\alpha p$, and $A_nm=-\alpha m$.  Consequently,
almost surely,
\begin{equation}
  \mathcal K_2(A_n,b)=\operatorname{span}\{p,m\}
  \label{eq:overview-krylov}
\end{equation}
is two-dimensional and $A_n$-invariant.  This identity is the main reduction:
the Ritz matrix $H=Q^TA_nQ$ is simply the restriction of $A_n$ to the span of
these two random spectral components, expressed in orthonormal coordinates.
In particular, $Q^T$ preserves the angle between $p$ and $m$.

The core technical point is to show that the particular components selected by
a random start inherit the exponentially small separation of the two
eigenspaces with high probability.  Let $s$ be the sine of the acute angle
between $p$ and $m$, and put $a=\delta u-v$.  The Gram determinant of the
scaled vectors $(a,0)$ and $(v,\delta v)$ gives
\begin{equation}
  s^2\leq \delta^2\left(1+
  \frac{\lVert u\rVert_2^2}{\lVert\delta u-v\rVert_2^2}\right).
  \label{eq:overview-angle}
\end{equation}
A Gaussian norm-ratio bound shows that
$\lVert u\rVert_2\leq2\lVert v\rVert_2$ except with probability at most
$2e^{-r/20}$.  On this event,
$\lVert\delta u-v\rVert_2\geq\lVert v\rVert_2/2$, so
$s<5\delta$.  The two unit eigenvectors of $H$ have this same angle, and the
elementary angle formula for a $2\times2$ eigenvector matrix yields
\begin{equation}
  \kappa_V(H)\geq\frac1s\geq\frac1{5\delta}
  =\frac{e^n}{5}.
  \label{eq:overview-conditioning}
\end{equation}
This proves the high-probability conditioning obstruction in
\Cref{thm:main}.

For the pseudospectral statement, large eigenvector condition number alone is
not used as a black box; the proof extracts a quantitative disk directly from
the same angle.  In an orthonormal basis adapted to one eigenvector, $H$ has
the real Schur form
\begin{equation}
  T=\begin{pmatrix}\alpha&t\\0&-\alpha\end{pmatrix},
  \qquad |t|=\delta\frac{\sqrt{1-s^2}}{s}.
  \label{eq:overview-schur}
\end{equation}
Although the eigenvalues $\pm\alpha$ are exponentially close, the estimate
$s<5\delta$ forces $|t|\geq1/10$.  For every $z\in\mathbb C$, the product of
the singular values of $zI-T$ is $|z^2-\alpha^2|$, whereas the larger singular
value is at least $|t|$.  Hence
\[
  \sigma_{\min}(zI-T)
  \leq\frac{|z^2-\alpha^2|}{|t|}
  \leq5\delta^2
  \qquad\text{whenever }|z|\leq\alpha.
\]
Thus the disk of radius $\alpha=\delta/2$ lies in
$\Lambda_{5\delta^2}(H)$ on the same high-probability event.

The disk has area $\pi\delta^2/4$.  Multiplying by the probability of the
norm-ratio event gives the expected-area lower bound in \Cref{thm:main}.
Finally, at $\varepsilon_n=5e^{-2n}$ this lower bound is exponentially larger
than $p(n)\varepsilon_n^\beta$ for every polynomial $p$ and every fixed
$\beta>1$.  The conditioning and pseudospectral conclusions therefore come
from one obstruction: orthonormal compression preserves the nearly coincident
spectral directions, and that same geometry becomes a constant Schur coupling
at the chosen exponential scale.

\section{Proof of the Main Theorem}
\label{sec:proof}

We first isolate the only probabilistic estimate needed in the construction.

\begin{lemma}[Gaussian norm-ratio concentration]
\label{lem:gaussian-ratio}
Let $u,v\in\mathbb R^r$ be independent standard Gaussian vectors.  Then
\[
  \Pr\{\lVert u\rVert_2\leq2\lVert v\rVert_2\}
  \geq1-2e^{-r/20}.
\]
\end{lemma}

\begin{proof}
For a standard Gaussian vector $g\in\mathbb R^r$ and
$X=\lVert g\rVert_2^2$, evaluation of the one-dimensional Gaussian integral
gives
\[
  \mathbb E e^{tX}=(1-2t)^{-r/2}\quad(0<t<1/2),
  \qquad
  \mathbb E e^{-tX}=(1+2t)^{-r/2}\quad(t>0).
\]
Markov's inequality, with $t=1/4$ for the first transform and $t=1/2$ for the
second, yields
\[
  \Pr\{X\geq2r\}
  \leq e^{-r/2}2^{r/2}\leq e^{-r/20},
  \qquad
  \Pr\{X\leq r/2\}
  \leq e^{r/4}2^{-r/2}\leq e^{-r/20}.
\]
Apply the upper-tail estimate to $\lVert u\rVert_2^2$ and the lower-tail
estimate to $\lVert v\rVert_2^2$.  Outside the union of the two exceptional
events,
$\lVert u\rVert_2^2\leq2r$ and $\lVert v\rVert_2^2\geq r/2$, which imply the
claimed norm ratio.
\end{proof}

We next define the input family and use the norm-ratio event to control the
angle between the two spectral components of the start.  This angle is exactly
the geometric quantity governing the eigenvector condition number after
orthonormal compression.

Let $n=2r\geq4$ be even, set
\[
  \delta=e^{-n},\qquad \alpha=\frac{\delta}{2},
\]
and define
\begin{equation}
  A_n=
  \begin{pmatrix}
    \alpha I_r&-I_r\\
    0&-\alpha I_r
  \end{pmatrix}.
  \label{eq:input-family}
\end{equation}
For a vector $b=(u,v)\in\mathbb R^r\oplus\mathbb R^r$, define its spectral
components
\begin{equation}
  p=(u-\delta^{-1}v,0),
  \qquad
  m=(\delta^{-1}v,v).
  \label{eq:spectral-components}
\end{equation}

\begin{lemma}[Krylov compression from nearly coincident eigenspaces]
\label{lem:counterexample-geometry}
Let $A_n$ be defined by \Cref{eq:input-family} and let
$b=(u,v)\sim N(0,I_n)$.  Then $A_n$ is diagonalizable,
$\lVert A_n\rVert_2\leq1+e^{-n}/2<3/2$, and
$\mathcal K_2(A_n,b)$ is two-dimensional almost surely.  On this event, let
$Q\in\mathbb R^{n\times2}$ be any orthonormal basis of the Krylov space.  Then
\[
  \Pr\!\left\{\kappa_V(Q^TA_nQ)\geq\frac1{5\delta}\right\}
  \geq1-2e^{-r/20}.
\]
\end{lemma}

\begin{proof}
The subspaces
\[
  E_+=\{(x,0):x\in\mathbb R^r\},
  \qquad
  E_-=\{(x,\delta x):x\in\mathbb R^r\}
\]
are complementary eigenspaces of $A_n$ with eigenvalues $\alpha$ and
$-\alpha$, respectively.  Indeed,
\[
  A_n=
  \begin{pmatrix}I_r&I_r\\0&\delta I_r\end{pmatrix}
  \begin{pmatrix}\alpha I_r&0\\0&-\alpha I_r\end{pmatrix}
  \begin{pmatrix}I_r&I_r\\0&\delta I_r\end{pmatrix}^{-1}.
\]
Thus $A_n$ is diagonalizable.  Also,
\[
  A_n=\operatorname{diag}(\alpha I_r,-\alpha I_r)
  +\begin{pmatrix}0&-I_r\\0&0\end{pmatrix},
\]
so $\lVert A_n\rVert_2\leq\alpha+1=1+e^{-n}/2<3/2$.

The vectors in \Cref{eq:spectral-components} satisfy
\[
  b=p+m,\qquad A_np=\alpha p,\qquad A_nm=-\alpha m.
\]
Consequently,
\begin{equation}
  \mathcal K_2(A_n,b)=\operatorname{span}\{p,m\}
  \label{eq:krylov-spectral-span}
\end{equation}
whenever $p$ and $m$ are nonzero.  For Gaussian $u,v$, this occurs almost
surely.  The vectors are then linearly independent because they belong to the
complementary eigenspaces $E_+$ and $E_-$, so the Krylov space is
two-dimensional and invariant under $A_n$.

Let $s$ be the sine of the acute angle between the lines spanned by $p$ and
$m$.  Scaling both vectors by $\delta$ does not alter this angle.  Put
$a=\delta u-v$; the scaled vectors are $(a,0)$ and $(v,\delta v)$.  Applying
\Cref{eq:gram-angle} gives
\begin{align*}
  s^2
  &=\frac{\lVert a\rVert_2^2(1+\delta^2)\lVert v\rVert_2^2-(a^Tv)^2}
  {\lVert a\rVert_2^2(1+\delta^2)\lVert v\rVert_2^2}\\
  &=\frac{\delta^2\bigl(
  \lVert u\rVert_2^2\lVert v\rVert_2^2-(u^Tv)^2
  +\lVert a\rVert_2^2\lVert v\rVert_2^2\bigr)}
  {\lVert a\rVert_2^2(1+\delta^2)\lVert v\rVert_2^2}\\
  &\leq\delta^2\left(1+
  \frac{\lVert u\rVert_2^2}{\lVert a\rVert_2^2}\right).
\end{align*}
On the event $\lVert u\rVert_2\leq2\lVert v\rVert_2$, the inequality
$\delta\leq1/4$ gives
\[
  \lVert a\rVert_2
  \geq\lVert v\rVert_2-\delta\lVert u\rVert_2
  \geq\frac12\lVert v\rVert_2.
\]
It follows that
\begin{equation}
  s\leq\sqrt{17}\,\delta<5\delta.
  \label{eq:small-angle}
\end{equation}

By \Cref{eq:invariant-compression}, the unit eigenvectors of
$H=Q^TA_nQ$ have the same acute angle as $p$ and $m$.  Therefore
\Cref{eq:angle-conditioning,eq:small-angle} give
\[
  \kappa_V(H)\geq\frac1s\geq\frac1{5\delta}.
\]
The probability estimate follows from
\Cref{lem:gaussian-ratio}.
\end{proof}

It remains to turn the same small-angle event into an area lower bound.  In an
orthonormal basis adapted to one Ritz eigenvector, the angle controls the
off-diagonal entry of a triangular form.  That entry supplies the singular
value estimate needed for a pseudospectral disk.

\begin{lemma}[Pseudospectral disk from Schur coupling]
\label{lem:pseudospectral-disk}
Under the assumptions of \Cref{lem:counterexample-geometry}, set
$\varepsilon_n=5\delta^2$.  Apart from the null event on which a spectral
component in \Cref{eq:spectral-components} vanishes, the event
$\lVert u\rVert_2\leq2\lVert v\rVert_2$ implies
\[
  \{z\in\mathbb C:|z|\leq\alpha\}
  \subseteq\Lambda_{\varepsilon_n}(Q^TA_nQ).
\]
Consequently,
\[
  \mathbb E\!\left[\operatorname{area}
  \Lambda_{\varepsilon_n}(Q^TA_nQ)\right]
  \geq(1-2e^{-r/20})\frac{\pi\delta^2}{4}.
\]
\end{lemma}

\begin{proof}
Write $H=Q^TA_nQ$.  Its eigenvalues are $\alpha$ and $-\alpha$.  Let $s$ and
$c=\sqrt{1-s^2}$ be the sine and cosine of the acute angle between its two unit
eigenvectors.  Choosing the first eigenvector and its orthogonal complement as
an orthonormal basis, the two eigenvector equations show that $H$ is
orthogonally similar to
\begin{equation}
  T=\begin{pmatrix}\alpha&t\\0&-\alpha\end{pmatrix},
  \qquad
  |t|=2\alpha\frac{c}{s}=\delta\frac{c}{s}.
  \label{eq:schur-form}
\end{equation}
On the stated event, \Cref{eq:small-angle} gives $s\leq5\delta$.  Because
$n\geq4$ implies $5\delta<1/2$, we have $c>1/2$, and hence
\begin{equation}
  |t|\geq\frac1{10}.
  \label{eq:schur-coupling}
\end{equation}

For every $z\in\mathbb C$, \Cref{eq:two-by-two-singular-values} and the entry
bound $\sigma_{\max}(zI-T)\geq|t|$ give
\[
  \sigma_{\min}(zI-T)
  \leq\frac{|z^2-\alpha^2|}{|t|}.
\]
If $|z|\leq\alpha$, then \Cref{eq:schur-coupling} makes the right-hand side at
most
\[
  10(|z|^2+\alpha^2)
  \leq20\alpha^2
  =5\delta^2
  =\varepsilon_n.
\]
Orthogonal similarity preserves singular values, proving the disk inclusion.
The disk has area $\pi\alpha^2=\pi\delta^2/4$.  Multiplying this deterministic
lower bound by the event probability in \Cref{lem:gaussian-ratio}, and using
nonnegativity of area on the complementary event, proves the expectation
bound.
\end{proof}

We now combine the geometric and pseudospectral estimates and verify both
uniform impossibility statements in \Cref{thm:main}.

\begin{proof}[Proof of \Cref{thm:main}]
Take the matrix $A_n$ in \Cref{eq:input-family}.  By
\Cref{lem:counterexample-geometry}, it is diagonalizable,
$\lVert A_n\rVert_2\leq1+e^{-n}/2<3/2$, and its two-step Gaussian Krylov space has
dimension two almost surely.  Since $r=n/2$ and $\delta=e^{-n}$, the same
lemma gives
\[
  \Pr\!\left\{\kappa_V(Q^TA_nQ)\geq\frac{e^n}{5}\right\}
  \geq1-2e^{-n/40}.
\]
For every fixed positive polynomial $p$, one has $e^n/5>p(n)$ for all sufficiently
large $n$.  Hence $\kappa_V(Q^TA_nQ)\leq p(n)$ fails with probability tending
to one exponentially fast along the even dimensions.

The expectation bound in the theorem is exactly
\Cref{lem:pseudospectral-disk} after substituting $r=n/2$ and
$\delta=e^{-n}$.  Finally, fix $\beta>1$ and a positive polynomial $p$.  At
$\varepsilon_n=5e^{-2n}$,
\begin{align*}
  \frac{\mathbb E[\operatorname{area}
  \Lambda_{\varepsilon_n}(Q^TA_nQ)]}
  {p(n)\varepsilon_n^\beta}
  &\geq
  \frac{(1-2e^{-n/40})\pi}
  {4\cdot5^\beta p(n)}e^{2(\beta-1)n}.
\end{align*}
This ratio tends to infinity along the even dimensions, so no uniform
$p(n)\varepsilon^\beta$ upper bound can hold.  This completes the proof.
\end{proof}

\bibliographystyle{alpha}
\bibliography{references}

\end{document}
